\section{Simulation Results}
\label{sec:simulation_results}

Report the simulation setup and measured performance. Use the same metric names
as Section~\ref{sec:spec_summary} so the final result is easy to compare with the
requirements.

\subsection{Simulation Environment}

Describe the simulator, process corner, temperature, supply voltage, load
condition, input condition, and any important testbench settings.

\begin{table}[H]
    \centering
    \caption{Simulation environment.}
    \label{tab:simulation_environment}
    \renewcommand{\arraystretch}{1.2}
    \begin{tabular}{l c}
        \toprule
        \textbf{Item} & \textbf{Setting} \\
        \midrule
        Process corner & TBD \\
        Temperature & TBD \\
        Supply voltage & TBD \\
        Load condition & TBD \\
        Input condition & TBD \\
        \bottomrule
    \end{tabular}
\end{table}

\subsection{Performance Summary}

Summarize the final measured metrics and explain whether each target is met.

\begin{table}[H]
    \centering
    \caption{Measured performance summary.}
    \label{tab:measured_performance}
    \renewcommand{\arraystretch}{1.2}
    \begin{tabularx}{0.9\textwidth}{
        >{\raggedright\arraybackslash}X
        >{\centering\arraybackslash}p{3.5cm}
        >{\centering\arraybackslash}p{3cm}
    }
        \toprule
        \textbf{Metric} & \textbf{Value} & \textbf{Status} \\
        \midrule
        Metric 1 & TBD & TBD \\
        Metric 2 & TBD & TBD \\
        Metric 3 & TBD & TBD \\
        \bottomrule
    \end{tabularx}
\end{table}

\subsection{Waveforms and Plots}

Include the key simulation waveforms, spectra, transient results, or corner
results. Each figure should state what is being measured and how the result
supports the specification.

\begin{figure}[H]
    \centering
    \fbox{\parbox[c][5cm][c]{0.85\textwidth}{
        \centering
        Replace this box with a simulation result plot.
    }}
    \caption{Key simulation result.}
    \label{fig:key_simulation_result}
\end{figure}
